% Shared by main.tex, main_spread.tex and the page checks. Compile from the project root.
\usepackage[paperwidth=182mm,paperheight=251mm,margin=0mm]{geometry}
\usepackage{fontspec}
\usepackage{graphicx}
\usepackage{tikz}
\usepackage{needspace}
\pagestyle{empty}
% paper: the plates are cleaned to this tone (tools/make_plates.py TONE)
\definecolor{paper}{RGB}{246,240,226}
\pagecolor{paper}

% body: the 版下 is a written 楷書 with katakana -> Klee One (OFL), vertical glyph forms
\newfontfamily\edbody{KleeOne-SemiBold.ttf}[Path=fonts/,RawFeature=+vert]
\newfontfamily\edruby{KleeOne-SemiBold.ttf}[Path=fonts/,RawFeature=+vert]
% preface (written 行草): 衡山毛筆フォント行書 (free, commercial use allowed)
\newfontfamily\edgyo{KouzanGyousho.ttf}[Path=fonts/,RawFeature=+vert]
% characters Klee One lacks: 全字庫正楷體
\newfontfamily\edfallback{TW-Kai.ttf}[Path=fonts/]
\newfontfamily\edfallbackB{TW-Kai-ExtB.ttf}[Path=fonts/]

\usepackage[hidelinks,bookmarksopen]{hyperref}

% Covers (front/cover.tex, front/back.tex)
\definecolor{slipink}{RGB}{48,40,30}
\definecolor{threadcol}{RGB}{236,230,214}
% \sewing{x}{dir}: 四つ目綴じ, holes at x mm, loops to the spine edge in direction dir
% (+1 right, -1 left); hole heights measured on the cover scans (17.5/88.5/162/233 mm from the top)
\newcommand\sewing[2]{%
  \def\edgex{\ifnum#2>0 182\else 0\fi}%
  \begin{scope}[threadcol,line width=0.75mm,line cap=round]
    \draw (#1,251) -- (#1,0);
    \foreach \y in {233.5,162.5,89,18}{
      \draw (#1,\y) -- (\edgex,\y);
      \fill[black!65] (#1,\y) circle (0.6);}
  \end{scope}}
\newcommand\edcover{\setpaper{182mm}\null% 書衣（表）: 柿渋色の紙に左上の題簽，右側に四つ目綴じ（右綴じの本）。
% Positions measured on the cover scans (qian_003, kun_001); textures: tools/make_covers.py.
% \edvol = qian | kun (set by the main file).
\begin{tikzpicture}[remember picture,overlay,shift={(current page.south west)},xshift=\edxshift,x=1mm,y=1mm]
  \node[anchor=south west,inner sep=0pt] at (0,0) {\includegraphics[width=182mm,height=251mm]{figures/cover_paper.jpg}};
  % 題簽 (slip 42.5 x 186.5 mm at 5 mm from the left, 2.5 mm from the top) with a soft shadow
  \fill[black,opacity=0.10] (5.6,61.4) rectangle (48.1,247.9);
  \node[anchor=north west,inner sep=0pt] at (5,248.5)
    {\includegraphics[width=42.5mm,height=186.5mm]{figures/slip_paper.jpg}};
  \draw[slipink,line width=0.5mm] (6.6,246.9) rectangle (45.9,63.6);
  \draw[slipink,line width=0.18mm] (7.7,245.8) rectangle (44.8,64.7);
  \node[anchor=north west,inner sep=0pt] at (5,248.5)
    {\includegraphics[width=42.5mm,height=186.5mm]{figures/slip_ink_\edvol.png}};
  \sewing{171}{1}
\end{tikzpicture}
\newpage}
\newcommand\edback{\setpaper{182mm}\null% 書衣（裏）: 同じ紙，四つ目綴じは左（表の裏返し）。下寄りに小さな刊記風の枠で本書の情報。
\begin{tikzpicture}[remember picture,overlay,shift={(current page.south west)},xshift=\edxshift,x=1mm,y=1mm]
  \node[anchor=south west,inner sep=0pt] at (0,0) {\includegraphics[width=182mm,height=251mm]{figures/cover_paper.jpg}};
  \sewing{11}{-1}
  \node[anchor=south,inner sep=0pt,text width=96mm,align=center,slipink] at (96,22) {%
    \apxrm\fontsize{8.4pt}{14pt}\selectfont
    \XeTeXlinebreaklocale "ja"\XeTeXlinebreakskip=0pt plus 1pt
    {\fontsize{12pt}{18pt}\selectfont 外蕃容貌圖画　\edvolname\par}\smallskip
    田川春道撰　倉田東岳画　安政二年（一八五五）刊\par
    底本　筑波大学附属図書館所蔵本（ヤ600-4）\par\smallskip
    {\fontsize{7.4pt}{12pt}\selectfont
    図版は筑波大学附属図書館デジタルコレクションの画像を改変して使用した。
    同館の利用条件により，調査・研究・教育・学習の目的に限る。\par
    翻字・組版　二〇二六年（Claude の支援による。未校閲）\par
    \ifdefined\ededition 中国語訳・英訳　二〇二六年（Claude による。専門家の校閲を経ていない）\par
    {\tren\fontsize{7.6pt}{11pt}\selectfont Chinese and English translations by Claude (Anthropic), not reviewed by a specialist.\par}\fi}};
  \ifdefined\ededition\def\edboxtop{61}\else\def\edboxtop{52}\fi
  \draw[slipink,line width=0.35mm] (46,18) rectangle (146,\edboxtop);
  \draw[slipink,line width=0.15mm] (47.2,19.2) rectangle (144.8,\edboxtop-1.2);
\end{tikzpicture}
\newpage}

% Back matter (凡例, 異体字一覧, 校記, 序 釈文): horizontal, modern setting with
% XeTeX's own CJK line breaking (no xeCJK: it would interfere with the page fonts)
\newfontfamily\apxrm{BIZ-UDMinchoM.ttc}[Path=fonts/]
\newfontfamily\apxsung{TW-Sung.ttf}[Path=fonts/]
\newfontfamily\apxsungB{TW-Sung-ExtB.ttf}[Path=fonts/]
\newfontfamily\apxgyo{KouzanGyousho.ttf}[Path=fonts/]
\newcommand\apxbig{\fontsize{14pt}{16pt}\selectfont}
\newenvironment{appendixpages}{%
  \newgeometry{top=26mm,bottom=26mm,left=24mm,right=24mm}%
  \apxrm\fontsize{10pt}{18pt}\selectfont
  \XeTeXlinebreaklocale "ja"\XeTeXlinebreakskip=0pt plus 1.5pt\emergencystretch=2em
  \setlength\parindent{1em}\setlength\parskip{3pt}\raggedbottom}{\restoregeometry}
\newcommand\apxsection[2]{\par\Needspace{10\baselineskip}\bigskip\hypertarget{#1}{}\pdfbookmark[0]{#2}{bm-#1}%
  {\centering\fontsize{15pt}{22pt}\selectfont #2\par}\medskip}
\newenvironment{vartable}{\par\medskip\centering
  \begin{tabular}{*{3}{c@{\enspace}c@{\enspace}r@{\hspace{3em}}}}}%
  {\end{tabular}\par\medskip}
% 校記: \jjhead{target}{page}{country}, then \jjitem{its name}{place in the text}{number}{text}
% per note (the place and the number at the column foot are written by tools/gen_pages.py)
\definecolor{linkcol}{RGB}{150,30,20}
\newcommand\jjhead[3]{\par\medskip\noindent\hyperlink{#1}{\textcolor{linkcol}{#2}}\quad #3\par}
\newcommand\jjitem[4]{\noindent\hangindent=2em\hangafter=1\hypertarget{#1}{}%
  \makebox[2em][l]{\hspace*{0.4em}\hyperlink{#2}{\textcolor{linkcol}{#3}}}#4\par}

% Translation pages of the trilingual edition (tools/gen_trans.py): a fixed-size block per page,
% Chinese in Noto Serif SC, English in EB Garamond (TeX Live), notes at the foot of the block
\newfontfamily\trzh{NotoSerifSC-Regular.ttf}[Path=fonts/,BoldFont=NotoSerifSC-SemiBold.ttf]
\newfontfamily\tren{EB Garamond}[Numbers=OldStyle]
\newfontfamily\trlat{EB Garamond}[Numbers=Lining]
\let\trcjk\trzh
\newif\iftrzh
\newenvironment{trpage}[3]{% {zh|en}{width}{height}
  \begin{minipage}[t][#3][t]{#2}\def\tmplang{#1}\def\tmpzh{zh}%
  \ifx\tmplang\tmpzh\trzhtrue\trzh\fontsize{13pt}{24.5pt}\selectfont
    \XeTeXlinebreaklocale "zh"\XeTeXlinebreakskip=0pt plus 1pt minus 0.1pt\parindent=2em
  \else\trzhfalse\tren\fontsize{13.5pt}{18.5pt}\selectfont
    \XeTeXlinebreaklocale ""\XeTeXlinebreakskip=0pt\parindent=1.5em\fi
  \parskip=0pt\emergencystretch=1.5em}{\end{minipage}}
% \trhead{link to the original page}{continent}{country}
\newcommand\trhead[3]{\noindent{\fontsize{9.5pt}{12pt}\selectfont\color{black!55}#2}\par\vspace{1.2mm}%
  \noindent\hyperlink{#1}{{\iftrzh\bfseries\fi\fontsize{18pt}{24pt}\selectfont #3}}\par\vspace{2.5mm}%
  \noindent{\color{black!45}\rule{\linewidth}{0.4pt}}\par\vspace{6mm}}
\newcommand\trtight{\iftrzh\fontsize{12pt}{22pt}\else\fontsize{12pt}{16.2pt}\fi\selectfont}
\newcommand\trmark[1]{\raisebox{0.65ex}{\fontsize{7pt}{7pt}\selectfont\trlat #1}}
\newcommand\trsmall[1]{{\iftrzh\fontsize{10.5pt}{24.5pt}\else\fontsize{11pt}{18.5pt}\fi\selectfont #1}}
\newcommand\trnotes{\par\vfill\noindent{\color{black!45}\rule{30mm}{0.4pt}}\par\vspace{1.5mm}%
  \iftrzh\fontsize{9pt}{14pt}\else\fontsize{9.5pt}{12.5pt}\fi\selectfont\parindent=0pt}
\newcommand\trnote[2]{\par\hangindent=1.3em\hangafter=1\noindent\makebox[1.3em][l]{\trlat #1}#2\par\vspace{0.6mm}}

\newdimen\edxshift \edxshift=0pt
% the page's destination p-<key> is set inside the page itself (works in both layouts)
\newcommand\edpage[1]{\null\input{pages/#1}\newpage}
% openings: one PDF page = right page (ura) + left page (next omote), as in the scans
\newcommand\setpaper[1]{\paperwidth=#1\relax\pdfpagewidth=#1\relax}
\newcommand\edspread[2]{%
  \setpaper{364mm}\null
  \if\relax\detokenize{#1}\relax\else{\edxshift=182mm\input{pages/#1}}\fi
  \if\relax\detokenize{#2}\relax\else{\edxshift=0pt\input{pages/#2}}\fi
  \newpage}
